% Points de contrôle de Flink et sortie transactionnelle : barrières dans le flux,
% instantané (positions, état, transaction), validation à chaque point de contrôle.
\documentclass[border=6pt]{standalone}
\usepackage{sd}
\begin{document}
\begin{tikzpicture}[x=1mm, y=1mm, note/.append style={font=\sffamily\small},
    op/.style={soft=#1, minimum width=24mm, minimum height=11mm, font=\sffamily\small, align=center},
    lab/.style={font=\sffamily\small, text=#1, align=center}, lab/.default=Dink]
  % ---- chaîne
  \node[op=Dslate] (k1) at (0,30) {Kafka « clics »\\{\scriptsize relisible}};
  \node[op=Dblue] (src) at (36,30) {source\\{\scriptsize position lue}};
  \node[op=Dblue] (win) at (72,30) {fenêtres\\{\scriptsize état : comptes}};
  \node[op=Dblue] (snk) at (108,30) {sortie\\{\scriptsize transaction ouverte}};
  \node[op=Dslate] (k2) at (144,30) {Kafka « comptes »\\{\scriptsize lu en read\_committed}};
  \foreach \a/\b in {k1/src, src/win, win/snk, snk/k2} \draw[arr=Dslate] (\a) -- (\b);
  \foreach \x in {50,86,122} {\draw[Dred, line width=1.6pt] (\x,26.5) -- (\x,33.5);}
  \node[lab=Dred] at (86,39) {barrière n : tout ce qui la précède appartient au point de contrôle n};
  \node[op=Damber, minimum width=70mm] (st) at (72,8) {point de contrôle n : positions, état des fenêtres, transaction à valider};
  \foreach \o in {src, win, snk} \draw[Damber, dash pattern=on 3pt off 2pt, line width=0.6pt, -{Stealth}] (\o.south) -- (\o.south |- st.north);
  % ---- chronologie des transactions
  \begin{scope}[yshift=-24mm]
    \node[note, anchor=east] at (-6,14) {exactly-once};
    \draw[ligne] (0,14) -- (170,14);
    \foreach \x/\t in {0/{T1}, 40/{T2}, 80/{T3}} {
      \draw[draw=Dteal, fill=DtealBg, line width=0.5pt] (\x+1,11.5) rectangle (\x+39,16.5);
      \node[lab=Dteal] at (\x+20,14) {\t\ validée};}
    \draw[draw=Dred, fill=DredBg, line width=0.5pt] (121,11.5) rectangle (137,16.5);
    \node[lab=Dred] at (129,14) {T4};
    \node[lab=Dred, font=\sffamily\bfseries] at (139,14) {×};
    \draw[draw=Dteal, fill=DtealBg, line width=0.5pt] (145,11.5) rectangle (169,16.5);
    \node[lab=Dteal] at (157,14) {T4 refaite};
    \node[lab=Dred, anchor=north] at (130,10) {panne : T4 annulée, jamais visible ;\\reprise depuis le point de contrôle 3};
    \node[note, anchor=east] at (-6,-6) {at-least-once};
    \draw[ligne] (0,-6) -- (170,-6);
    \draw[draw=Dslate, fill=DslateBg, line width=0.5pt] (1,-8.5) rectangle (137,-3.5);
    \node[lab=Dslate] at (69,-6) {résultats écrits et visibles aussitôt};
    \draw[draw=Dred, fill=DredBg, line width=0.5pt] (145,-8.5) rectangle (169,-3.5);
    \node[lab=Dred] at (157,-6) {réécrits};
    \node[lab=Dred, anchor=north] at (150,-10) {ce qui suivait le point 3 sort deux fois\\(2 650 à 3 000 résultats en double, mesuré)};
  \end{scope}
\end{tikzpicture}
\end{document}
